\documentclass[aspectratio=169,10pt]{beamer}
% Shared preamble for the Algorithms and Programming Concepts slides (Beamer, Metropolis)
\usetheme[progressbar=frametitle,numbering=fraction,background=light,block=fill,sectionpage=none]{metropolis}
\usepackage[T1]{fontenc}
\usepackage{FiraSans}
\usepackage[scaled=.85]{FiraMono}
\usepackage{amsmath,amssymb}
\usefonttheme[onlymath]{serif}
\usepackage{booktabs,array,tabularx}
\usepackage{tikz,pgfplots}
\usepackage{xcolor}
\usepackage{listings}
\usepackage{adjustbox}
\usepackage{fontawesome5}
\hyphenpenalty=10000\exhyphenpenalty=10000
\newenvironment{sfig}[1][0.80]{\begin{adjustbox}{max width=\linewidth,max totalheight=#1\textheight,center}}{\end{adjustbox}}
\usetikzlibrary{arrows.meta,positioning,shapes.geometric,shapes.misc,fit,calc,backgrounds,decorations.pathreplacing,patterns,matrix,chains}
\pgfplotsset{compat=1.17}

\definecolor{tdInk}{HTML}{1D2330}
\definecolor{tdMuted}{HTML}{5D6678}
\definecolor{tdLine}{HTML}{D9D5CA}
\definecolor{tdPaper}{HTML}{FBFAF7}
\definecolor{tdBlue}{HTML}{2457D6}
\definecolor{tdGreen}{HTML}{17855A}
\definecolor{tdAmber}{HTML}{B7791F}
\definecolor{tdRed}{HTML}{C8362F}
\definecolor{tdPurple}{HTML}{7A4FD1}
\setbeamercolor{normal text}{fg=tdInk,bg=tdPaper}
\setbeamercolor{alerted text}{fg=tdRed}
\setbeamercolor{example text}{fg=tdGreen}
\setbeamercolor{frametitle}{bg=tdInk,fg=white}
\setbeamercolor{progress bar}{fg=tdBlue,bg=tdLine}
\setbeamercolor{title separator}{fg=tdBlue}
\setbeamercolor{block title}{bg=tdBlue!12,fg=tdInk}
\setbeamercolor{block body}{bg=tdBlue!5,fg=tdInk}
\setbeamercolor{block title alerted}{bg=tdRed!12,fg=tdRed!80!black}
\setbeamercolor{block body alerted}{bg=tdRed!5,fg=tdInk}
\setbeamercolor{block title example}{bg=tdGreen!14,fg=tdGreen!70!black}
\setbeamercolor{block body example}{bg=tdGreen!6,fg=tdInk}
\setbeamerfont{frametitle}{size=\large,series=\bfseries}
\makeatletter\setlength{\metropolis@frametitle@padding}{1.15ex}\makeatother
\setbeamertemplate{navigation symbols}{}
\setbeamerfont{block title}{size=\small}

% ---- code listings (Python 3.14) ----
\lstdefinestyle{py}{language=Python,
  basicstyle=\ttfamily\fontsize{8.4}{10.1}\selectfont\color{tdInk},
  keywordstyle=\color{tdBlue}\bfseries, commentstyle=\color{tdMuted}\itshape, stringstyle=\color{tdGreen},
  morekeywords={match,case,type,self,True,False,None},
  emph={print,len,range,sorted,min,max,sum,enumerate,zip,heappush,heappop,cache},emphstyle=\color{tdPurple},
  showstringspaces=false, columns=fullflexible, keepspaces=true, frame=single, rulecolor=\color{tdLine},
  backgroundcolor=\color{white}, xleftmargin=2pt, xrightmargin=2pt, aboveskip=2pt, belowskip=0pt, framesep=3pt}
\lstdefinestyle{out}{basicstyle=\ttfamily\fontsize{8}{9.6}\selectfont\color{tdInk}, columns=fullflexible, keepspaces=true,
  frame=single, rulecolor=\color{tdGreen!60}, backgroundcolor=\color{tdGreen!6}, xleftmargin=2pt, xrightmargin=2pt, aboveskip=2pt, belowskip=0pt, framesep=3pt}
% \pycode[size]{snippet} shows snippets/<snippet>.py ; \pyout[size]{snippet} shows snippets/<snippet>.out
\newcommand{\pycode}[2][\relax]{{#1\lstinputlisting[style=py]{snippets/#2.py}}}
\newcommand{\pyout}[2][\relax]{{#1\lstinputlisting[style=out]{snippets/#2.out}}}

% ---- tikz styles ----
\tikzset{>={Stealth[length=2.2mm]},
  card/.style={rounded corners=4pt,draw=#1!70!black,fill=#1!10,align=center,inner sep=5pt,font=\small}, card/.default=tdBlue,
  pill/.style={rounded corners=7pt,draw=none,fill=#1,text=white,inner sep=3pt,font=\scriptsize\bfseries,align=center}, pill/.default=tdBlue,
  arr/.style={->,thick,draw=tdMuted},
  cell/.style={draw=tdInk!70,minimum width=7.5mm,minimum height=7.5mm,inner sep=0pt,font=\small\ttfamily,fill=white},
  tn/.style={circle,draw=tdInk!80,fill=white,minimum size=7mm,inner sep=0pt,font=\small\bfseries}}
\newcommand{\key}[1]{\textbf{\color{tdBlue}#1}}
\newcommand{\bad}[1]{\textbf{\color{tdRed}#1}}
\newcommand{\good}[1]{\textbf{\color{tdGreen!70!black}#1}}
\newcommand{\apcfooter}[1]{\setbeamertemplate{footline}{\hfill\tiny\color{tdMuted}Algorithms and Programming Concepts \,|\, #1 \,|\, BUV\hspace{1.6em}\insertframenumber/\inserttotalframenumber\hspace{0.8em}\vspace{4pt}}}
\newcommand{\apctitle}[2]{\title{#1}\subtitle{#2}\author{Dr.\ Quang-Vinh Dang\\ {\small Senior Lecturer (Assoc.\ Prof)}}\institute{British University Vietnam}\date{Algorithms and Programming Concepts \,|\, Year 1 \,|\, Python 3.14}}

\apcfooter{Optimization}
\apctitle{Optimization}{Finding the best answer among too many candidates}

\begin{document}
\begin{frame}[plain]\titlepage\end{frame}

\begin{frame}{What is optimization?}
\begin{columns}[T]
\begin{column}{0.5\textwidth}
{\small
Most real problems are not ``find \emph{an} answer'' but ``find the \textbf{best} answer'':
\begin{itemize}\setlength\itemsep{3pt}
\item the shortest delivery route;
\item the most valuable load for a truck;
\item the model parameters with the smallest error.
\end{itemize}
\vspace{1mm}
Every optimization problem has three parts:
\begin{enumerate}\setlength\itemsep{2pt}
\item \textbf{decision variables}: what we choose;
\item \textbf{constraints}: what is allowed;
\item an \textbf{objective}: the number to maximise or minimise.
\end{enumerate}}
\end{column}
\begin{column}{0.46\textwidth}
\begin{block}{After this lecture you can}
{\small
\begin{itemize}\setlength\itemsep{2pt}
\item write the three parts of a problem;
\item solve \textbf{knapsack} by brute force, greedy and DP;
\item explain \textbf{local} versus \textbf{global} optimum;
\item run \textbf{hill climbing} and \textbf{gradient descent};
\item see why \textbf{TSP} is hard and why heuristics are useful.
\end{itemize}}
\end{block}
\end{column}
\end{columns}
\end{frame}

\begin{frame}{The knapsack problem}
\begin{columns}[T]
\begin{column}{0.5\textwidth}
\pycode{knap_items}
\vspace{2mm}
{\small
A bag holds weight $50$. Each item can be taken or left. Which items give the largest total value?

\vspace{2mm}
\begin{tabular}{@{}lccc@{}}\toprule item & weight & value & value/weight\\\midrule A & 10 & 60 & 6.0\\ B & 20 & 100 & 5.0\\ C & 30 & 120 & 4.0\\\bottomrule\end{tabular}}
\end{column}
\begin{column}{0.46\textwidth}
\begin{sfig}[0.7]
\begin{tikzpicture}
\foreach \i/\n/\w/\v/\c in {0/A/10/60/tdGreen,1/B/20/100/tdBlue,2/C/30/120/tdAmber}{
 \node[card,fill=\c!15,minimum width=1.6cm,minimum height=0.9cm+\w*0.03cm] at (\i*2.1,0) {\textbf{\n}\\\scriptsize w=\w, v=\v};}
\node[draw=tdInk!60,rounded corners,minimum width=5.8cm,minimum height=0.8cm] at (2.1,-2) {bag: capacity 50};
\end{tikzpicture}
\end{sfig}
{\small
With $n$ items there are $2^n$ ways to choose. Here: $2^3=8$ subsets, easy to check.}
\end{column}
\end{columns}
\end{frame}

\begin{frame}{Attempt 1: brute force}
\begin{columns}[T]
\begin{column}{0.54\textwidth}
\pycode{knap_brute}
\end{column}
\begin{column}{0.42\textwidth}
{\small
Try \textbf{every} subset (\texttt{combinations}), keep the best one that fits. Always correct.

\vspace{2mm}
Result for our items:}
\pyout{knap_demo}
{\small
\vspace{2mm}
The best choice is B + C with value 220. But there are $2^n$ subsets: what happens when $n$ grows?}
\end{column}
\end{columns}
\end{frame}

\begin{frame}{Brute force does not scale}
\begin{columns}[T]
\begin{column}{0.5\textwidth}
\pyout{knap_growth}
\vspace{2mm}
{\small
Each 4 extra items multiplies the subsets by 16, and the time with it. For 18 items brute force needs \textbf{0.18 s}. For 40 items it would need days. The DP column barely moves: we come to it in a few slides.}
\end{column}
\begin{column}{0.46\textwidth}
\begin{sfig}[0.78]
\begin{tikzpicture}
\begin{axis}[width=8cm,height=5.4cm,ymode=log,xlabel={items $n$},ylabel={subsets $2^n$},axis lines=left,tick label style={font=\scriptsize},label style={font=\small}]
\addplot[tdRed,very thick,domain=1:40,samples=40] {2^x};
\end{axis}
\end{tikzpicture}
\end{sfig}
\end{column}
\end{columns}
\end{frame}

\begin{frame}{Attempt 2: greedy, the tempting shortcut}
\begin{columns}[T]
\begin{column}{0.52\textwidth}
\pycode{knap_greedy}
\end{column}
\begin{column}{0.44\textwidth}
{\small
Idea: take items with the best \textbf{value per kilo} first, while they fit. Fast: sorting, $O(n\log n)$.

\vspace{2mm}
Greedy takes A (ratio 6), then B (5). C no longer fits. Total $160$.}
\pyout{knap_demo}
{\small
\vspace{1mm}
\bad{Greedy is not optimal here:} B + C gives 220. Greedy never looks back, and the best bag needs a worse-ratio item.

\emph{Greedy is a heuristic, not a guarantee.}}
\end{column}
\end{columns}
\end{frame}

\begin{frame}{Attempt 3: dynamic programming}
\begin{columns}[T]
\begin{column}{0.54\textwidth}
\pycode{knap_dp}
\end{column}
\begin{column}{0.42\textwidth}
{\small
\texttt{best[c]} = largest value that fits in capacity \texttt{c}. For each item we ask for every \texttt{c}: ``is it better to take this item, or skip it?''

\vspace{2mm}
\texttt{max(best[c], best[c - w] + v)}

\vspace{2mm}
The loop runs backwards so each item is used at most once. Work: $O(n\times\text{capacity})$, \textbf{polynomial} in the number of items.

\vspace{2mm}
It gives 220, like brute force, in microseconds.}
\end{column}
\end{columns}
\end{frame}

\begin{frame}{Three methods compared}
\renewcommand{\arraystretch}{1.4}\small
\begin{tabular}{@{}>{\bfseries}p{26mm}>{\raggedright}p{38mm}>{\raggedright}p{36mm}>{\raggedright\arraybackslash}p{40mm}@{}}
\toprule & Brute force & Greedy & Dynamic programming\\\midrule
Optimal? & \good{yes} & \bad{not always} & \good{yes}\\
Time & $O(2^n)$ & $O(n\log n)$ & $O(n\cdot W)$\\
Idea & try everything & best ratio first & reuse sub-answers\\
Use when & $n$ tiny & a good-enough answer is fine & sub-problems overlap\\\bottomrule
\end{tabular}

\vspace{4mm}
{\small $W$ is the capacity. The DP time is polynomial in the \emph{value} of $W$, not in its number of digits; this is why knapsack is still a hard problem in theory.}
\end{frame}

\begin{frame}{Continuous problems: climbing a hill}
\begin{columns}[T]
\begin{column}{0.5\textwidth}
\begin{sfig}[0.85]
\begin{tikzpicture}
\begin{axis}[width=8.4cm,height=5.8cm,xmin=0,xmax=10,ymin=-4,ymax=10,xlabel={$x$},ylabel={$f(x)=x\sin x$},axis lines=left,tick label style={font=\scriptsize},label style={font=\small},clip=false]
\addplot[tdBlue,very thick,domain=0:10,samples=120] {x*sin(deg(x))};
\addplot[only marks,mark=*,mark size=2.6pt,tdAmber] coordinates {(2.03,1.82)};
\addplot[only marks,mark=*,mark size=2.6pt,tdGreen] coordinates {(7.98,7.92)};
\node[font=\scriptsize,tdAmber!80!black,anchor=south] at (axis cs:2.03,2.3) {local peak};
\node[font=\scriptsize,tdGreen!60!black,anchor=south] at (axis cs:7.98,8.4) {global peak};
\end{axis}
\end{tikzpicture}
\end{sfig}
\end{column}
\begin{column}{0.46\textwidth}
{\small
Now $x$ is a number in $[0,10]$ and we want the largest $f(x)$. We cannot try all real numbers.

\vspace{2mm}
\key{Hill climbing}: stand somewhere, look one small step left and right, move to the higher side, and stop when neither side is higher.

\vspace{2mm}
The curve has two peaks. The \textbf{global} optimum is the highest point; a \textbf{local} optimum is only higher than its neighbours.}
\end{column}
\end{columns}
\end{frame}

\begin{frame}{Hill climbing in Python}
\begin{columns}[T]
\begin{column}{0.54\textwidth}
\pycode{hill}
\end{column}
\begin{column}{0.42\textwidth}
\pyout{hill_demo}
{\small
\vspace{2mm}
Starting at $x=1$ we climb the small hill and \bad{get stuck at a local peak} ($f=1.82$). Starting at $x=6$ we find the global peak ($f=7.92$).

\vspace{2mm}
The result depends on where we start.}
\end{column}
\end{columns}
\end{frame}

\begin{frame}{Escaping local peaks: random restarts}
\begin{columns}[T]
\begin{column}{0.54\textwidth}
\pycode{restarts}
\end{column}
\begin{column}{0.42\textwidth}
{\small
Climb from many random starting points and keep the best peak found. With 10 restarts the global peak was found ($x=7.98$, $f=7.92$, see the previous slide).

\vspace{2mm}
No guarantee, but it is simple and often works. Better-known variants add controlled randomness, such as \textbf{simulated annealing} and \textbf{genetic algorithms}.}
\end{column}
\end{columns}
\end{frame}

\begin{frame}{Gradient descent: use the slope}
\begin{columns}[T]
\begin{column}{0.54\textwidth}
\pycode{gradient}
\vspace{2mm}
\pyout{gd_demo}
\end{column}
\begin{column}{0.42\textwidth}
{\small
To \emph{minimise} $f(x)=(x-3)^2+1$ we move against the slope:
\[x \leftarrow x - \eta\, f'(x)\]
$\eta$ (\texttt{lr}) is the \textbf{learning rate}, the step size. This single idea trains neural networks.

\vspace{2mm}
\good{lr = 0.1:} smooth approach to $x=3$.\\
\bad{lr = 1.0:} jumps between 10 and $-4$ forever.\\
\bad{lr = 1.1:} each step overshoots more; the value explodes.

\vspace{2mm}
\emph{Too small is slow, too large diverges.}}
\end{column}
\end{columns}
\end{frame}

\begin{frame}{The learning rate in a picture}
\begin{columns}[T]
\begin{column}{0.58\textwidth}
\begin{sfig}[0.85]
\begin{tikzpicture}
\begin{axis}[width=9.6cm,height=6cm,xmin=0,xmax=12,ymin=-20,ymax=20,xlabel={step},ylabel={$x$},axis lines=left,legend pos=north east,legend style={font=\scriptsize},tick label style={font=\scriptsize},label style={font=\small},clip=true]
\addplot[tdGreen,very thick,mark=*,mark size=1.4pt] table[x=step,y=lr01]{snippets/gd_data.out};
\addplot[tdAmber,very thick,mark=*,mark size=1.4pt] table[x=step,y=lr10]{snippets/gd_data.out};
\addplot[tdRed,very thick,mark=*,mark size=1.4pt] table[x=step,y=lr11]{snippets/gd_data.out};
\addplot[dashed,tdInk!60] coordinates {(0,3) (12,3)};
\legend{lr=0.1,lr=1.0,lr=1.1,minimum $x=3$}
\end{axis}
\end{tikzpicture}
\end{sfig}
\end{column}
\begin{column}{0.4\textwidth}
{\small
Data from the real run (12 steps, start at $x=10$).

\vspace{2mm}
Green walks to the minimum. Amber bounces between two values. Red swings wider each time.

\vspace{2mm}
In machine learning, picking the learning rate is one of the most important settings of the training.}
\end{column}
\end{columns}
\end{frame}

\begin{frame}{The travelling salesman problem (TSP)}
\begin{columns}[T]
\begin{column}{0.58\textwidth}
\pycode{tsp_nn}
\end{column}
\begin{column}{0.4\textwidth}
{\small
Visit every city once and return home, with the shortest total distance.

\vspace{2mm}
\textbf{Nearest neighbour} (greedy): always go to the closest unvisited city. It is fast, $O(n^2)$, but short-sighted.

\vspace{2mm}
\texttt{tour\_length} adds the distance of each leg, including the return to the start.}
\end{column}
\end{columns}
\end{frame}

\begin{frame}{TSP: the exact answer by brute force}
\begin{columns}[T]
\begin{column}{0.52\textwidth}
\pycode{tsp_brute}
\vspace{2mm}
{\small
Try all $(n-1)!$ orders (\texttt{permutations}) and keep the shortest. City 0 is fixed as the start.}
\end{column}
\begin{column}{0.44\textwidth}
{\small On 8 cities:}
\pyout{tsp_demo}
{\small
\vspace{2mm}
The greedy tour is 25\% longer than the best. Is exactness worth it? It depends on how many cities we have.}
\end{column}
\end{columns}
\end{frame}

\begin{frame}{Why TSP is hard}
\begin{columns}[T]
\begin{column}{0.5\textwidth}
\pyout{tsp_growth}
\vspace{2mm}
{\small
Eight cities: 5\,040 tours, instant. Twenty cities: $1.2\times10^{17}$ tours. At a billion tours per second that is about \textbf{four years}.}
\end{column}
\begin{column}{0.46\textwidth}
{\small
Nearest neighbour was 25\% longer than the best tour, but it needs no search at all, so it scales to thousands of cities.

\vspace{2mm}
No one knows a fast exact algorithm for TSP; it belongs to the family of \textbf{NP-hard} problems. In practice we use:
\begin{itemize}\setlength\itemsep{2pt}
\item \textbf{heuristics} that are fast and usually good;
\item \textbf{local search} that improves a tour (e.g.\ 2-opt);
\item \textbf{approximation} with a guaranteed quality.\end{itemize}}
\end{column}
\end{columns}
\end{frame}

\begin{frame}{Choosing a strategy}
\renewcommand{\arraystretch}{1.35}\small
\begin{tabular}{@{}>{\raggedright\bfseries}p{38mm}>{\raggedright}p{50mm}>{\raggedright\arraybackslash}p{58mm}@{}}
\toprule Situation & Method & Warning\\\midrule
Few candidates & brute force & $2^n$ and $n!$ explode quickly\\
Overlapping sub-problems & dynamic programming & you must find the right state\\
Smooth function, slope known & gradient descent & learning rate; local minima\\
Rough landscape, no slope & hill climbing + restarts & may stop at a local peak\\
Huge hard problem (TSP) & heuristics, local search & no guarantee of the best answer\\\bottomrule
\end{tabular}

\vspace{4mm}
{\small Searching, sorting, graphs and DP from the earlier lectures are the building blocks: optimization is where they are used together.}
\end{frame}

\begin{frame}{Recap and practice}
\begin{columns}[T]
\begin{column}{0.48\textwidth}
\begin{block}{Take-aways}
{\small
\begin{enumerate}\setlength\itemsep{3pt}
\item An optimization problem = \key{variables, constraints, objective}.
\item \key{Brute force} is exact but exponential.
\item \key{Greedy} is fast but may be wrong (knapsack).
\item \key{DP} is exact and polynomial when sub-problems overlap.
\item \key{Local search} can stop at a local optimum; restarts help.
\item \key{Gradient descent} needs a good learning rate.
\item \key{TSP}: we accept good heuristic answers.
\end{enumerate}}
\end{block}
\end{column}
\begin{column}{0.48\textwidth}
\begin{block}{Practice}
{\small
\begin{enumerate}\setlength\itemsep{3pt}
\item Find a 3-item knapsack where greedy gives the worst possible answer.
\item Make \texttt{knapsack\_dp} return the \emph{chosen items}, not only the value.
\item Run \texttt{hill\_climb} from 20 different starts. How many reach the global peak?
\item Try \texttt{lr = 0.5} and \texttt{lr = 0.9}. Which converges faster than 0.1?
\item Implement 2-opt: reverse a part of the tour while it gets shorter. Compare with nearest neighbour.
\end{enumerate}}
\end{block}
\end{column}
\end{columns}
\end{frame}
\end{document}
